\exercisesfor{1}

\begin{enumerate}
\def\labelenumi{\arabic{enumi}.}
\item
  Let \(\mathfrak{g}\) be a Lie algebra and \(\mathfrak{a}\) and \(\mathfrak{b}\) ideals (resp. characteristic ideals) of \(\mathfrak{g}\) . Then the set of \(x \in \mathfrak{g}\) such that \([x, \mathfrak{b}] \subset \mathfrak{a}\) is an ideal (resp. a characteristic ideal) of \(\mathfrak{g}\) called the transporter of \(\mathfrak{b}\) into \(\mathfrak{a}\) . Show that \(\mathcal{C}_{p+1}\mathfrak{g}\) is the transporter of \(\mathfrak{g}\) into \(\mathcal{C}_p\mathfrak{g}\) .
\item
  If \(\mathfrak{g}\) is an \(n\) -dimensional Lie algebra over a field \(\mathbf{K}\) and the centre of \(\mathfrak{g}\) is of dimension \(\geqslant n - 1\) , \(\mathfrak{g}\) is commutative.
\item
  Let M be a not necessarily associative algebra over a field K, \((e_{i})_{i\in I}\) a basis of the vector space M and \(c_{ijk}\) the constants of structure of M relative to the basis \((e_{i})\) . For M to be a Lie algebra, it is necessary and sufficient that the \(c_{ijk}\) satisfy the following conditions: (a) \(c_{iik}=0\) ; (b) \(c_{ijk}=-c_{jik}\) ; (c) \(\sum_{r\in I}(c_{ijr}c_{rkl}+c_{jkr}c_{ril}+c_{kir}c_{rjl})=0\) for all i,j,k,l in I.
\item
  \begin{enumerate}
  \def\labelenumii{(\alph{enumii})}
  \tightlist
  \item
    Suppose that 2 is invertible in K. Let \(\mathfrak{a}\) be a Lie algebra over K. On the K-module \(\mathfrak{a}' = \mathfrak{a} \times \mathfrak{a}\) we define a bracket by the formula:
  \end{enumerate}
\end{enumerate}

\[
[ (x _ {1}, x _ {2}), (y _ {1}, y _ {2}) ] = ([ x _ {1}, y _ {1} ] + [ x _ {2}, y _ {2} ], [ x _ {1}, y _ {2} ] + [ x _ {2}, y _ {1} ]).
\]

Then \(\mathfrak{a}'\) is a Lie K-algebra and the mappings \(x\mapsto \frac{1}{2} (x,x)\) , \(x\mapsto \frac{1}{2} (x, - x)\) are isomorphisms of \(\mathfrak{a}\) onto ideals \(\mathfrak{b}\) , \(\mathfrak{c}\) of \(\mathfrak{a}'\) of which \(\mathfrak{a}'\) is the direct sum. (Consider the quadratic extension \(\mathrm{K}'\) of \(\mathbf{K}\) with basis 1 and \(k\) where \(k^2 = 1\) . Form \(\mathfrak{a}_{(\mathrm{K}',1)}\) and then restrict the ring of scalars to \(\mathbf{K}\) . Then observe that \(\frac{1}{2} (1 + k)\) , \(\frac{1}{2}(1 - k)\) are idempotents of \(\mathrm{K}'\) and that \((1 + k)(1 - k) = 0\) .)

\begin{enumerate}
\def\labelenumi{(\alph{enumi})}
\setcounter{enumi}{1}
\tightlist
\item
  Let \(a\) be a real Lie algebra, \(\mathfrak{g} = a_{(C)}\) and \(b\) the real Lie algebra derived from \(\mathfrak{g}\) by restricting the field of scalars to \(\mathbf{R}\) . Show that \(b_{(C)}\) is isomorphic to \(a_{(C)}'\) , where \(a'\) is the algebra introduced in (a). (Consider \(\mathrm{V} = \mathbf{C} \otimes_{\mathbf{R}} \mathbf{C}\) as a vector space over \(\mathbf{C}\) by \((z \otimes z')z'' = z \otimes z'z''\) . Observe that \(\mathrm{V}\) is the complexification of the real vector subspace of \(\mathrm{V}\) generated by \(1 \otimes 1\) and \(i \otimes i\) and that this subspace is identified with the quadratic extension of \(\mathbf{R}\) with basis 1 and \(k\) where \(k^2 = 1\) .)
\end{enumerate}

\begin{enumerate}
\def\labelenumi{\arabic{enumi}.}
\setcounter{enumi}{4}
\item
  Let \(\mathrm{V}\) and \(\mathrm{W}\) be vector spaces of dimensions \(n + 1\) and \(n\) over a field \(\mathbf{K}\) . Show that \(\mathfrak{gl}(\mathbf{W})\) is isomorphic to a subalgebra of \(\mathfrak{sl}(\mathbf{V})\) . (W can be assumed to be a hyperplane of \(\mathrm{V}\) . Let \(e \in \mathbb{V}\) , \(e \notin \mathbb{W}\) . For \(x \in \mathfrak{gl}(\mathbf{W})\) , let \(y(x)\) be the element of \(\mathfrak{gl}(\mathbf{V})\) which extends \(x\) and is such that \(y(e) = -\mathrm{Tr}(x).e\) . Show that the mapping \(x \mapsto y(x)\) is an isomorphism of \(\mathfrak{gl}(\mathbf{W})\) onto a subalgebra of \(\mathfrak{sl}(\mathbf{V})\) .
\item
  For a Lie algebra g to be associative, it is necessary and sufficient that \(\mathcal{D}\mathfrak{g}\) be contained in the centre of \(\mathfrak{g}\) .
\end{enumerate}

¶ 7. Let A be a ring (possibly without unit element) such that the relation \(2a = 0\) in A implies \(a = 0\) . Let U be a subring of A and an ideal of the Lie Z-algebra associated with A.

\begin{enumerate}
\def\labelenumi{(\alph{enumi})}
\tightlist
\item
  If \(x, y\) belong to U, then \((xy - yx)\mathrm{A} \subset \mathrm{U}\) (write:
\end{enumerate}

\[
(x y - y x) s = x (y s) - (y s) x - y (x s - s x))
\]

and \(\mathrm{A}(xy - yx)\mathrm{A} \subset \mathrm{U}\) (write:

\[
r (x y - y x) s = (x y - y x) s r + r [ (x y - y x) s ] - [ (x y - y x) s ] r).
\]

\begin{enumerate}
\def\labelenumi{(\alph{enumi})}
\setcounter{enumi}{1}
\item
  Suppose that U is commutative. Let \(x \in \mathbf{U}\) , \(s \in \mathbf{A}\) and \(y = xs - sx\) . Show that \((yt)^3 = 0\) for all \(t \in \mathbf{A}\) . (The element \(x\) commutes with \(xs - sx\) and \(xs^2 - s^2x\) ; deduce that \(2(xs - sx)^2 = 0\) , whence \(y^2 = 0\) . Similarly \((yt - ty)^2 = 0\) for all \(t \in \mathbf{A}\) .)
\item
  Suppose henceforth that the only two-sided ideals of A are \(\{0\}\) and A and that for every non-zero element y of A there exists \(t \in A\) such that yt is not nilpotent. Show that, if \(U \neq A\) , U is contained in the centre of A. (Show, using \(\langle a\rangle\) , that U is commutative and then use \(\langle b\rangle\) .)
\item
  Let V be an ideal of the Lie Z-algebra associated with A. Then either V is contained in the centre of A or \(V \supset [A, A]\) . (Apply (c) to the set U of \(t \in A\) such that \([t, A] \subset V\) ; use the identity \([xy, z] = [x, yz] + [y, zx]\) .)
\end{enumerate}

\begin{enumerate}
\def\labelenumi{\arabic{enumi}.}
\setcounter{enumi}{7}
\tightlist
\item
  If \(x_{1}, x_{2}, x_{3}, x_{4}\) are 4 elements of a Lie algebra, then
\end{enumerate}

\[
\begin{array}{r l} {[ [ [ x _ {1}, x _ {2} ], x _ {3} ], x _ {4} ] + [ [ [ x _ {2}, x _ {1} ], x _ {4} ], x _ {3} ]} & \\ & \quad + [ [ [ x _ {3}, x _ {4} ], x _ {1} ], x _ {2} ] + [ [ [ x _ {4}, x _ {3} ], x _ {2} ], x _ {1} ] = 0. \end{array}
\]

\begin{enumerate}
\def\labelenumi{\arabic{enumi}.}
\setcounter{enumi}{8}
\tightlist
\item
  Let \(\mathfrak{g}\) be a Lie algebra in which \([[x,y],y] = 0\) for all \(x\) and \(y\) .
\end{enumerate}

\begin{enumerate}
\def\labelenumi{(\alph{enumi})}
\item
  Show that \(3[[x,y],z] = 0\) for all \(x,y,z\) in \(\mathfrak{g}\) . (Observe that \([[x,y],z]\) is an alternating trilinear function of \((x,y,z)\) and apply the Jacobi identity.)
\item
  Show similarly, using (a) and Exercise 8, that \(\left[[[x,y],z],t\right]=0\) for all \(x,y,z,t\) in \(\mathfrak{g}\) .
\end{enumerate}

\begin{enumerate}
\def\labelenumi{\arabic{enumi}.}
\setcounter{enumi}{9}
\tightlist
\item
  Let L be an associative algebra and g the associated Lie algebra. Every derivation of L is a derivation of g, but the converse is not true. (Consider the identity mapping of a commutative associative algebra.)
\end{enumerate}
% semantic-fragments: legacy-input-seam
\relax{}
\begin{enumerate}
\def\labelenumi{\arabic{enumi}.}
\setcounter{enumi}{10}
\item
  Let \(\mathfrak{g}\) be a Lie algebra and \(\mathfrak{h}\) an ideal of \(\mathfrak{g}\) such that \(\mathcal{D}\mathfrak{h} = \mathfrak{h}\) . Show that \(\mathfrak{h}\) is a characteristic ideal of \(\mathfrak{g}\) .
\item
  Let g be a Lie algebra.
\end{enumerate}

\begin{enumerate}
\def\labelenumi{(\alph{enumi})}
\item
  For a derivation D of g to commute with all the inner derivations of g, it is necessary and sufficient that D map g into its centre.
\item
  Suppose henceforth that \(\mathfrak{g}\) has centre zero, so that \(\mathfrak{g}\) is identified with an ideal of the Lie algebra \(\mathfrak{D}\) of all derivations of \(\mathfrak{g}\) . Show that the centralizer of \(\mathfrak{g}\) in \(\mathfrak{D}\) is zero (use (a)). In particular, the centre of \(\mathfrak{D}\) is zero.
\item
  Show that a derivation \(\Delta\) of \(\mathfrak{D}\) such that \(\Delta(\mathfrak{g}) = \{0\}\) is zero. \(([\Delta(\mathfrak{D}),\mathfrak{g}]\subset \Delta([\mathfrak{D},\mathfrak{g}]) + [\mathfrak{D},\Delta(\mathfrak{g})]\subset \Delta(\mathfrak{g}) = \{0\}\) and hence \(\Delta(\mathfrak{D}) = \{0\}\) by (b).)
\item
  Show that, if further \(\mathfrak{g} = \mathcal{D}\mathfrak{g}\) , every derivation \(\Delta\) of \(\mathfrak{D}\) is inner (use (c) and Exercise 11).
\end{enumerate}

\begin{enumerate}
\def\labelenumi{\arabic{enumi}.}
\setcounter{enumi}{12}
\tightlist
\item
  Let \(\mathfrak{g}\) be a Lie algebra and \(\mathfrak{a}\) and \(\mathfrak{b}\) two submodules of \(\mathfrak{g}\) . For every integer \(i \geqslant 0\) we define the submodules \(\mathfrak{m}_i = \mathfrak{m}_i(\mathfrak{a},\mathfrak{b})\) and \(\mathfrak{n}_i = \mathfrak{n}_i(\mathfrak{a},\mathfrak{b})\) as follows: \(\mathfrak{m}_1 = \mathfrak{b}, \mathfrak{m}_{i+1} = [\mathfrak{m}_i,\mathfrak{b}], \mathfrak{n}_1 = \mathfrak{a}, \mathfrak{n}_{i+1} = [\mathfrak{n}_i,\mathfrak{b}]\) . Show that
\end{enumerate}

\[
[ a, m _ {i} ] \subset n _ {i + 1}
\]

for \(i = 1,2,\ldots\) (argue by induction on \(i\) , observing that

\[
\mathfrak {m} _ {i} ([ a, b ], b) = \mathfrak {m} _ {i} (a, b)
\]

and \(n_{i}([a,b],b)=n_{i+1}(a,b)\) .

¶ 14. Let \(\mathfrak{a}\) be a Lie algebra and \(\mathfrak{b}\) a subalgebra of \(\mathfrak{a}\) . A composition series joining \(\mathfrak{a}\) to \(\mathfrak{b}\) is a decreasing sequence \((\mathfrak{a}_i)_{0\leqslant i\leqslant n}\) of subalgebras of \(\mathfrak{a}\) such that \(\mathfrak{a}_0 = \mathfrak{a}\) , \(\mathfrak{a}_n = \mathfrak{b}\) , \(\mathfrak{a}_{i+1}\) is an ideal of \(\mathfrak{a}_i\) for \(0\leqslant i < n\) . \(\mathfrak{b}\) is called a subinvariant subalgebra of \(\mathfrak{a}\) if there exists a composition series joining \(\mathfrak{a}\) to \(\mathfrak{b}\) .

\begin{enumerate}
\def\labelenumi{(\alph{enumi})}
\item
  Let \(\mathfrak{b}\) be a subinvariant subalgebra of \(\mathfrak{a}\) and \((\mathfrak{a}_i)_{0\leqslant i\leqslant n}\) a composition series joining \(\mathfrak{a}\) to \(\mathfrak{b}\) . Deduce from Exercise 13 that \([\mathcal{C}^{n+p}\mathfrak{b},\mathfrak{a}]\subset\mathcal{C}^{p}\mathfrak{b}\) by observing that \([\mathfrak{a}_i,\mathfrak{b}]\subset\mathfrak{a}_{i+1}\) for \(1\leqslant i<n\) .
\item
  Deduce from (a) that the intersection \(\mathcal{C}^{\infty}\mathfrak{b}\) of the \(\mathcal{C}^{p}\mathfrak{b}\) ( \(p = 0,1,2,\ldots\) ) is an ideal of \(\mathfrak{a}\) .
\item
  Deduce from (b) that a subinvariant subalgebra \(\mathfrak{c}\) of \(\mathfrak{a}\) such that \(\mathfrak{c} = \mathcal{D}\mathfrak{c}\) is an ideal of \(\mathfrak{a}\) .
\item
  When K is a field and \(\dim_{K}a<+\infty\) , show that the intersection \(D^{\infty}b\) of the \(D^{p}b\left(p=0,1,2,\ldots\right)\) is an ideal of a. (Show that this intersection is a subinvariant subalgebra of a and apply (c).)
\end{enumerate}

¶ 15. (a) Let a be a Lie algebra, b a subalgebra of a, c an ideal of b and z the centralizer of c in a. Then {[}b, b{]} ⊂ z. Deduce that b + z is a subalgebra of a in which z is an ideal.

\begin{enumerate}
\def\labelenumi{(\alph{enumi})}
\setcounter{enumi}{1}
\item
  Let g be a Lie algebra, a a subinvariant subalgebra of g and b an ideal of a. If the centralizer of b in a and the centralizer of a in g are zero, the centralizer \(\mathfrak{z}\) of \(\mathfrak{b}\) in \(\mathfrak{g}\) is zero. (Let \(\mathfrak{b} = \mathfrak{a} + \mathfrak{z}\) , which is a subalgebra by \((a)\) ; \(\mathfrak{a}\) is subinvariant in \(\mathfrak{b}\) ; if \(\mathfrak{z} \neq \{0\}\) , then \(\mathfrak{a} \neq \mathfrak{h}\) , hence \(\mathfrak{a}\) is an ideal in \(\mathfrak{a}'\) , with \(\mathfrak{a} \neq \mathfrak{a}' \subset \mathfrak{b}\) ; let \(a' \in \mathfrak{a}'\) , \(a' \notin \mathfrak{a}\) and \(a' = a + z\) , \(a \in \mathfrak{a}\) , \(z \in \mathfrak{z}\) , \(z \neq 0\) ; show that \([z, \mathfrak{a}]\) is contained in \(\mathfrak{a}\) and commutes with \(\mathfrak{b}\) , whence \([z, \mathfrak{a}] = \{0\}\) giving a contradiction.)
\item
  Let \(\mathfrak{g}\) be a Lie algebra with centre zero, \(\mathfrak{D}_1\) the Lie algebra of derivations of \(\mathfrak{g}\) , \(\mathfrak{D}_2\) the Lie algebra of derivations of \(\mathfrak{D}_1, \ldots, \mathfrak{g}\) is identified with an ideal of \(\mathfrak{D}_1\) , \(\mathfrak{D}_1\) with an ideal of \(\mathfrak{D}_2, \ldots\) (cf.~Exercise 12 (b)). Using (b) and Exercise 12 (b), show that the centralizer of \(\mathfrak{g}\) in \(\mathfrak{D}_i\) is zero for all \(i\) . (For a sequel to this exercise, cf.~§ 4, Exercise 15.)
\end{enumerate}

\begin{enumerate}
\def\labelenumi{\arabic{enumi}.}
\setcounter{enumi}{15}
\tightlist
\item
  Let g be a Lie algebra and D the Lie algebra of derivations of g. The identity mapping of D defines a semi-direct product h of D by g called the holomorph of g. g is identified with an ideal and D with a subalgebra of h.
\end{enumerate}

\begin{enumerate}
\def\labelenumi{(\alph{enumi})}
\item
  Show that the centralizer \(\mathfrak{g}^*\) of \(\mathfrak{g}\) in \(\mathfrak{h}\) is the set of \(\operatorname{ad}_{\mathfrak{g}}x - x(x\in \mathfrak{g})\) .
\item
  For every element \(x + d\) of \(\mathfrak{h}\) ( \(x \in \mathfrak{g}, d \in \mathfrak{D}\) ), we write
\end{enumerate}

\[
\theta (x + d) = \mathrm{ad} _ {\mathfrak {g}} x - x + d.
\]

Show that \(\theta\) is an automorphism of \(\mathfrak{h}\) of order 2 which maps \(\mathfrak{g}\) onto \(\mathfrak{g}^*\) , so that \(\mathfrak{h}\) can be considered as the holomorph of \(\mathfrak{g}^*\) .

\begin{enumerate}
\def\labelenumi{(\alph{enumi})}
\setcounter{enumi}{2}
\item
  For every ideal \(\mathfrak{k} \neq \{0\}\) of \(\mathfrak{h}\) , \(\mathfrak{k} \cap (\mathfrak{g} + \mathfrak{g}^*) \neq \{0\}\) . (If \(\mathfrak{k} \cap \mathfrak{g} = \{0\}\) , \(\mathfrak{k} \subset \mathfrak{g}^*\) .)
\item
  If g is commutative, every ideal \(t \neq \{0\}\) of h contains g. If K is a field and \(\dim_{K} g < +\infty\) , deduce that, if \(t \neq g\) , h cannot be the holomorph of t (use (b)).
\end{enumerate}

\begin{enumerate}
\def\labelenumi{\arabic{enumi}.}
\setcounter{enumi}{16}
\tightlist
\item
  Suppose that K is a field. Let T be an indeterminate and \(L = K((T))\) ; L has canonically a valuation (order of the formal power series) which makes it a complete valued field.
\end{enumerate}

\begin{enumerate}
\def\labelenumi{(\alph{enumi})}
\item
  Let \(\mathrm{V}\) be a finite-dimensional vector space over \(\mathbf{K}\) . Then \(\mathrm{V}_{(\mathbb{L})}\) has canonically a complete metrizable topological vector space structure (Topological Vector Spaces, Chapter I, § 2, no. 3, Theorem 2). Let \((x_{p})_{p\in \mathbb{Z}}\) be a family of elements of \(\mathrm{V}\) such that \(x_{p} = 0\) for \(p\) less than some rational integer. Then the series \(\sum_{-\infty}^{+\infty}x_p\mathrm{T}^p\) converges in \(\mathrm{V}_{(\mathrm{L})}\) and every element of \(\mathrm{V}_{(\mathrm{L})}\) can be expressed uniquely in this way. (Take a basis of V.)
\item
  The vector space \(\mathcal{L}_{\mathrm{L}}(\mathrm{V}_{(\mathrm{L})}) = \mathcal{L}_{\mathrm{K}}(\mathrm{V})_{(\mathrm{L})}\) has a canonical complete metrizable topological vector space structure. Every endomorphism of \(\mathrm{V}_{(\mathrm{L})}\) can be expressed uniquely in the form \(\sum_{-\infty}^{+\infty}u_p\mathrm{T}^p\) , where \(u_{p}\in \mathcal{L}_{\mathrm{K}}(\mathrm{V})\) and the \(u_{p}\) are zero for \(p\) less than some rational integer.
\item
  If \(\mathbf{K}\) is of characteristic 0 and \(u\in \mathcal{L}_{\mathrm{K}}(\mathrm{V})\) , we write
\end{enumerate}

\[
e ^ {\mathrm{Tu}} = 1 + \frac {1}{1 !} u \mathrm{T} + \frac {1}{2 !} u ^ {2} \mathrm{T} ^ {2} + \dots \in \mathcal {L} _ {\mathrm{L}} (\mathrm{V} _ {(\mathrm{L})}).
\]

If \(x \in V\) is such that \(ux = 0\) , then \(e^{Tu}.x = x\) .

\begin{enumerate}
\def\labelenumi{(\alph{enumi})}
\setcounter{enumi}{3}
\item
  If W is another finite-dimensional vector space over K and \(v \in \mathcal{L}_{\mathrm{K}}(\mathrm{W})\) , then \(e^{\mathrm{T}(u \otimes 1 + 1 \otimes v)} = e^{\mathrm{T}u} \otimes e^{\mathrm{T}v}\) .
\item
  Deduce from (c) and (d) that, if V has a (not necessarily associative) algebra structure and u is a derivation of V, then \(e^{Tu}\) is an automorphism of the algebra \(V_{(L)}\) .
\end{enumerate}

¶ 18. Let H be a complex Hilbert space, Y and Z continuous Hermitian operators on H and B = {[}Z, Y{]}. Suppose that Y and B are permutable.

\begin{enumerate}
\def\labelenumi{(\alph{enumi})}
\item
  Show that \([Z, Y^n] = nBY^{n-1}\) (argue by induction on \(n\) ).
\item
  Let \(f\) be a continuously differentiable function of a real variable. Deduce from (a) that \([Z, f(Y)] = Bf'(Y)\) .
\item
  Deduce from (b) that \(B = 0\) . (Show that, if \(B \neq 0\) , \(f\) can be chosen such that \(\| f(Y) \| \leqslant 1\) and \(\| Bf'(Y) \|\) is arbitrarily large.)
\item
  Let g be the Lie algebra of continuous operators T on H such that \(T^{*} = -T\) . Deduce from (c) that the conditions \(Y \in g\) , \(Z \in g\) , \([[Z, Y], Y] = 0\) imply \([Z, Y] = 0\) .
\end{enumerate}

\begin{enumerate}
\def\labelenumi{\arabic{enumi}.}
\setcounter{enumi}{18}
\tightlist
\item
  \begin{enumerate}
  \def\labelenumii{(\alph{enumii})}
  \tightlist
  \item
    Let \(\mathbf{L}\) be an associative algebra over \(\mathbf{K}\) . Given \(k\) elements \(x_{1},\ldots ,x_{k}\) of \(\mathbf{L}\) , we write:
  \end{enumerate}
\end{enumerate}

\[
f _ {k} (x _ {1}, \dots , x _ {k}) = \sum_ {\sigma \in \mathfrak {S}_{k}} x _ {\sigma (1)} \dots x _ {\sigma (k)}.
\]

In the algebra \(L \otimes_{K} K[T_{1}, \ldots, T_{k}]\) , where the \(T_{i}\) are indeterminates, \(f_{k}(x_{1}, \ldots, x_{k})\) is the coefficient of \(T_{1} \ldots T_{k}\) in \((x_{1}T_{1} + \cdots + x_{k}T_{k})^{k}\) ; it is also the coefficient of \(T_{1} \ldots T_{k}\) in:

\[
\sum_ {j = 0} ^ {k - 1} \left(x _ {1} \mathrm{T} _ {1} + \dots + x _ {k - 1} \mathrm{T} _ {k - 1}\right) ^ {j} x _ {k} \mathrm{T} _ {k} \left(x _ {1} \mathrm{T} _ {1} + \dots + x _ {k - 1} \mathrm{T} _ {k - 1}\right) ^ {k - j - 1}.
\]

Given two elements \(x, y\) in \(\mathbf{L}\) , we define \(g_{i,k-i}(x,y)\) as the coefficient of \(\mathrm{T}_1^i\mathrm{T}_2^{k - i}\) in \((x\mathrm{T}_1 + y\mathrm{T}_2)^k\) . Show that

\[
i! (k - i)! g _ {i, k - i} (x, y) = f _ {k} \left(t _ {1}, \dots , t _ {k}\right)
\]

where \(t_h = x\) for \(h \leqslant i\) and \(t_h = y\) for \(h \geqslant i + 1\) .

\begin{enumerate}
\def\labelenumi{(\alph{enumi})}
\setcounter{enumi}{1}
\tightlist
\item
  If L is considered as a Lie algebra over K, then in \(\mathcal{L}_{\mathrm{K}}(\mathrm{L})\) :
\end{enumerate}

\[
(\operatorname{ad} x) ^ {n} = \sum_ {i = 0} ^ {n} (- 1) ^ {n - i} \binom {n} {i} \mathrm{L} _ {x} ^ {i} R _ {x} ^ {n - i},
\]

where \(\mathbf{L}_x\) (resp. \(R_{x}\) ) is the multiplication \(y\mapsto xy\) (resp. \(y\mapsto yx\) ).

Deduce that, if K is of characteristic p (p prime), then:

\begin{enumerate}
\def\labelenumi{(\arabic{enumi})}
\tightlist
\item
\end{enumerate}

\[
(\operatorname{ad} x) ^ {p}. y = (\operatorname{ad} x ^ {p}). y\tag{2}
\]

\[
(\operatorname{ad} x) ^ {p - 1}. y = \sum_ {i = 0} ^ {p - 1} x ^ {i} y x ^ {p - 1 - i}.
\]

Deduce from (2) that:

\[
f _ {p - 1} (\operatorname{ad} x _ {1}, \dots , \operatorname{ad} x _ {p - 1}). y = f _ {p} (x _ {1}, \dots , x _ {p - 1}, y)
\]

and:

\[
g _ {i, p - 1 - i} (\mathrm{ad} x, \mathrm{ad} y). y = (p - i) g _ {i, p - i} (x, y).
\]

Conclude that for any two elements x, y of L:

\[
(x + y) ^ {p} = x ^ {p} + y ^ {p} + \Lambda_ {p} (x, y)\tag{3}
\]

where:

\[
\Lambda_ {p} (x, y) = \sum_ {i = 1} ^ {p - 1} (p - i) ^ {- 1} g _ {i, p - 1 - i} (\mathrm{ad} x, \mathrm{ad} y). y
\]

belongs to the Lie subalgebra of L generated by x and y (``Jacobson's formulae'').

\begin{enumerate}
\def\labelenumi{\arabic{enumi}.}
\setcounter{enumi}{19}
\tightlist
\item
  Let g be a Lie algebra over a ring K such that \(pK = \{0\}\) (p prime). A mapping \(x \mapsto x^{[p]}\) of g into itself is called a p-mapping if it satisfies the relations
\end{enumerate}

\[
\begin{array}{r l} & {\operatorname{ad} x ^ {[ p ]} = (\operatorname{ad} x) ^ {p}} \\ & {(\lambda x) ^ {[ p ]} = \lambda^ {p} x ^ {[ p ]}} \\ & {(x + y) ^ {[ p ]} = x ^ {[ p ]} + y ^ {[ p ]} + \Lambda_ {p} (x, y) \quad (\mathrm{cf.Exercise19(b)})} \end{array}
\]

  for \(x, y\) in \(\mathfrak{g}\) and \(\lambda \in \mathbf{K}\) . A Lie \(p\) -algebra over \(\mathbf{K}\) is a set with the structure defined by giving it a Lie algebra structure over \(\mathbf{K}\) and a \(p\) -mapping. Every associative algebra over \(\mathbf{K}\) is a Lie \(p\) -algebra with \([x, y] = xy - yx\) and \(x^{[p]} = x^p\) (Exercise 19). A mapping \(u\) of a Lie \(p\) -algebra \(\mathfrak{g}\) into a Lie \(p\) -algebra \(\mathfrak{g}'\) is called a \(p\) -homomorphism if \(u\) is a Lie algebra homomorphism and \(u(x^{[p]}) = (u(x))^{[p]}\) . Show that in a Lie \(p\) -algebra \(\mathfrak{g}\) every \(p\) -mapping is of the form \(x \mapsto x^{[p]} + f(x)\) , where \(f\) is a mapping of \(\mathfrak{g}\) into its centre, which is semi-linear for the endomorphism \(\lambda \mapsto \lambda^p\) of \(\mathbf{K}\) . More generally, if \(u\) is a homomorphism of \(\mathfrak{g}\) onto a Lie \(p\) -algebra \(\mathfrak{g}', u(x^{[p]}) - (u(x))^{[p]}\) belongs to the centre of \(\mathfrak{g}'\) .

\begin{enumerate}
\def\labelenumi{\arabic{enumi}.}
\setcounter{enumi}{20}
\tightlist
\item
  \begin{enumerate}
  \def\labelenumii{(\alph{enumii})}
  \tightlist
  \item
    Let \(\mathbf{K}\) be a field of characteristic \(p > 0\) and \(\mathbf{L}\) a not necessarily associative algebra over \(\mathbf{K}\) . Show that the derivations of \(\mathbf{L}\) form a Lie \(p\) -algebra with the \(p\) -mapping \(\mathbf{D} \mapsto \mathbf{D}^p\) .
  \end{enumerate}
\end{enumerate}

\begin{enumerate}
\def\labelenumi{(\alph{enumi})}
\setcounter{enumi}{1}
\item
  Let \(\mathbf{L}\) be the algebra \(\mathrm{K}[\mathrm{X}] / (\mathrm{X}^p)\) and \(\mathfrak{w}\) the \(p\) -algebra of derivations of \(\mathbf{L}\) . Let \(\mathrm{D}_i\) ( \(0 \leqslant i \leqslant p - 1\) ) denote the derivation of \(\mathbf{L}\) such that \(\mathrm{D}_i(\mathbf{X}) = \mathbf{X}^i\) . Show that the \(\mathrm{D}_i\) form a basis of \(\mathfrak{w}\) over \(\mathbf{K}\) and that \([\mathrm{D}_i, \mathrm{D}_j] = (j - i)\mathrm{D}_{i + j - 1}\) if \(i + j \leqslant p\) , \([\mathrm{D}_i, \mathrm{D}_j] = 0\) if \(i + j > p\) , and \(\mathrm{D}_i^p = 0\) , except for \(i = 1\) , in which case \(\mathrm{D}_1^p = \mathrm{D}_1\) .
\item
  Show that, if \(p \geqslant 3\) , the Lie algebra \(\mathfrak{w}\) only admits the ideals \(\{0\}\) and \(\mathfrak{w}\) . (Using the multiplication table of the \(D_i\) , show that every non-zero ideal of \(\mathfrak{w}\) contains a non-zero multiple of \(D_{p-1}\) .)
\item
  Let \(V_{k}\) be the subspace of w with basis the \(D_{i}\) of index \(\geqslant k\) . Show that, if \(p \geqslant 5\) , \(V_k\) is identical with the set of \(Z \in w\) whose centralizer is of dimension \(\geqslant k + 1\) . Deduce that, if \(p \geqslant 5\) , the automorphism group of \(w\) is solvable.
\end{enumerate}

\begin{enumerate}
\def\labelenumi{\arabic{enumi}.}
\setcounter{enumi}{21}
\tightlist
\item
  Let g be a Lie p-group over a ring K of characteristic p (p prime). Let \(x \mapsto x^{[p]}\) denote the p-mapping. For every subset E of g, let \(E^{p}\) denote the submodule of g generated by the \(x^{p}\) , where \(x \in E\) . An ideal \(a \subset g\) is called a p-ideal if \(a^{p} \subset a\) .
\end{enumerate}

\begin{enumerate}
\def\labelenumi{(\alph{enumi})}
\item
  For an ideal \(\mathfrak{a} \subset \mathfrak{g}\) to be a \(p\) -ideal, it is necessary and sufficient that it be the kernel of a \(p\) -homomorphism (Exercise 20).
\item
  The sum of two \(p\) -ideals is a \(p\) -ideal. The sum of a \(p\) -ideal and a \(p\) -subalgebra is a \(p\) -subalgebra.
\item
  Let \(\mathfrak{h}\) be a Lie subalgebra of \(\mathfrak{g}\) . The smallest Lie \(p\) -subalgebra containing \(\mathfrak{h}\) is \(\mathfrak{h} + \mathfrak{h}^p + \mathfrak{h}^{p^2} + \cdots = \overline{\mathfrak{h}}\) ; if we write \(\mathfrak{h}_i = \mathfrak{h} + \mathfrak{h}^p + \cdots + \mathfrak{h}^{p^i}\) , \(\mathfrak{h}_i\) is an ideal of \(\overline{\mathfrak{h}}\) and \(\overline{\mathfrak{h}}/\mathfrak{h}\) is commutative. If \(\mathfrak{h}\) is an ideal of \(\mathfrak{g}\) , so are the \(\mathfrak{h}_i\) and \(\overline{\mathfrak{h}}\) and the latter is the smallest \(p\) -ideal containing \(\mathfrak{h}\) .
\item
  The normalizer and centralizer of an arbitrary submodule of g are p-subalgebras of g.
\end{enumerate}

\begin{enumerate}
\def\labelenumi{\arabic{enumi}.}
\setcounter{enumi}{22}
\tightlist
\item
  Let \(\mathfrak{g}\) be a finite-dimensional commutative Lie \(p\) -algebra over a perfect field \(\mathbf{K}\) of characteristic \(p > 0\) .
\end{enumerate}

\begin{enumerate}
\def\labelenumi{(\alph{enumi})}
\item
  Show that g is in a unique way the direct sum of two p-subalgebras h, t such that on h the p-mapping is bijective and on t it is nilpotent (consider the iterations of the p-mapping on g; cf.~Algebra, Chapter VIII, § 2, no. 2, Lemma 2 and Chapter VII, § 5, Exercise 20 (a)). h is called the p-core of g.
\item
  Suppose that K is algebraically closed. Show that there exists a basis \((e_{i})\) of \(\mathfrak{h}\) such that \(e_{i}^{p}=e_{i}\) for all i (consider a p-subalgebra of h generated by a single element and show that it contains some x such that \(x^{p}=x\neq0\) ; then argue by induction on the dimension of \(\mathfrak{h}\)). Deduce that the p-subalgebras of h are finite in number and are in one-to-one correspondence with the vector subspaces of the vector space over the prime field \(F_{p}\) generated by the \(e_{i}\) .
\item
  Show that there exists a basis \((f_{ij})\) of \(\mathfrak{t}\) ( \(1 \leqslant i \leqslant k\) , \(1 \leqslant j \leqslant s_i\) for all \(i\) ) where the sequence of \(s_i\) is decreasing and which is such that \(f_{1j}^p = 0\) for \(1 \leqslant j \leqslant s_1\) , \(f_{ij}^p = f_{i-1,j}\) for \(2 \leqslant i \leqslant k\) , \(1 \leqslant j \leqslant s_i\) (method of Algebra, Chapter VII, § 5, Exercise 20 (b)).
\end{enumerate}

¶ 24. Let K be a (commutative) field of (arbitrary) characteristic p and \(X_{i}, Y_{i}, Z_{i}\) 3n distinct indeterminates; to abbreviate we shall denote the systems \((\mathrm{X}_{i})_{1 \leq i \leq n}, (\mathrm{Y}_{i})_{1 \leq i \leq n}, (\mathrm{Z}_{i})_{1 \leq i \leq n}\) by x, y, z respectively. Let \(A_{x,y,z}\) be the algebra of formal power series in the 3n indeterminates \(X_{i}, Y_{i}, Z_{i}\) with coefficients in K; we shall denote by \(A_{x,y}\) (resp. \(A_{x}\) ) the subalgebra of \(A_{x,y,z}\) consisting of the power series not containing z (resp. y, z). An element of \(A_{x,y,z}\) (resp. \(A_{x,y}, A_{x}\) ) will be denoted by \(u(\mathbf{x}, \mathbf{y}, \mathbf{z})\) (resp. \(u(\mathbf{x}, \mathbf{y}), u(\mathbf{x})\) );

 a system \((u_{i}(\mathbf{x},\mathbf{y},\mathbf{z}))_{1\leqslant i\leqslant n}\) of n elements of \(A_{x,y,z}\) will be denoted by \(u(\mathbf{x},\mathbf{y},\mathbf{z})\) ; analogous notation will be used for formal power series containing only one or two of the three systems x, y, z. If \(u(\mathbf{x},\mathbf{y},\mathbf{z})\in A_{x,y,z}\) and \(\mathbf{f}=(f_{i}),\mathbf{g}=(g_{i}),\mathbf{h}=(h_{i})\) are three systems of n elements of \(A_{x,y,z}\) which are formal power series with no constant term, we denote by \(u(\mathbf{f}(\mathbf{x},\mathbf{y},\mathbf{z}),\mathbf{g}(\mathbf{x},\mathbf{y},\mathbf{z}),\mathbf{h}(\mathbf{x},\mathbf{y},\mathbf{z}))\) the formal power series obtained by substituting \(f_{i}\) for \(X_{i}, g_{i}\) for \(Y_{i}, h_{i}\) for \(Z_{i}\) ( \(1\leqslant i\leqslant n\) ) in u. For every system \(\alpha=(\alpha_{1},\ldots,\alpha_{n})\) of n integers \(\geqslant0\) , let \(x^{\alpha}\) denote the monomial \(X_{1}^{\alpha_{1}}\ldots X_{n}^{\alpha_{n}}\) ; \(y^{\alpha}\) and \(z^{\alpha}\) are defined similarly. Let \(\varepsilon_{i}\) denote the system \(\alpha\) for which \(\alpha_{j}=0\) if \(j\neq i, \alpha_{i}=1\) . We denote by e the system \((0,\ldots,0)\) of n elements of \(A_{x,y,z}\) .

\begin{enumerate}
\def\labelenumi{(\alph{enumi})}
\item
  A formal group law over K (or, by an abuse of language, a formal group over K) of dimension n is a system \(G = \mathbf{f}(\mathbf{x}, \mathbf{y})\) of n elements of \(A_{x,y}\) with the following properties: (1) \(\mathbf{f}(\mathbf{x}, \mathbf{f}(\mathbf{y}, \mathbf{z})) = \mathbf{f}(\mathbf{f}(\mathbf{x}, \mathbf{y}), \mathbf{z})\) ; (2) \(\mathbf{f}(\mathbf{e}, \mathbf{y}) = \mathbf{y}\) , \(\mathbf{f}(\mathbf{x}, \mathbf{e}) = \mathbf{x}\) . Show that necessarily \(f_{i}(\mathbf{x}, \mathbf{y}) = \mathrm{X}_{i} + \mathrm{Y}_{i} + g_{i}(\mathbf{x}, \mathbf{y})\) , where \(g_{i}\) contains only monomials of degree \(\geqslant 2\) , each of which contains at least one \(X_{j}\) and at least one \(Y_{j}\) . Show that there exists one and only one system \(\mathbf{h}(\mathbf{x})\) of n elements of \(A_{x}\) such that \(\mathbf{f}(\mathbf{x}, \mathbf{h}(\mathbf{x})) = \mathbf{f}(\mathbf{h}(\mathbf{x}), \mathbf{x}) = \mathbf{e}\) (Algebra, Chapter IV, § 5, no. 9, Proposition 10). G is called commutative if \(\mathbf{f}(\mathbf{y}, \mathbf{x}) = \mathbf{f}(\mathbf{x}, \mathbf{y})\) .
\item
  For all \(u \in \mathbf{A}_{\mathbf{x}}\) , let \(L_{\mathbf{y}}u\) denote the element \(u(\mathbf{f}(\mathbf{y},\mathbf{x}))\) of \(\mathbf{A}_{\mathbf{x},\mathbf{y}}\) . A derivation D of \(\mathbf{A}_{\mathbf{x}}\) can be canonically extended to a derivation of \(\mathbf{A}_{\mathbf{x},\mathbf{y}}\) (also denoted D by an abuse of notation) by the condition that \(\mathrm{D}(\mathrm{Y}_i) = 0\) for all \(i\) (Algebra, Chapter IV, § 5, no. 8, Proposition 6). D is said to be left invariant for the formal group G in question if \(L_{\mathbf{y}}D = DL_{\mathbf{y}}\) . We denote by \(\mathrm{D}^{(\mathrm{e})}\) the linear mapping of \(\mathbf{A}_{\mathbf{x},\mathbf{y}}\) into \(\mathbf{A}_{\mathbf{y}}\) defined by \(\mathrm{D}^{(\mathrm{e})}(u) = (\mathrm{D}u)(\mathbf{e},\mathbf{y})\) ; it maps \(\mathbf{A}_{\mathbf{x}}\) into K and is determined by its restriction \(\mathbf{A}_{\mathbf{x}}\) . Show that, for D to be left invariant, it is necessary and sufficient that for all \(v \in \mathbf{A}_{\mathbf{x}}\) , \(\mathrm{D}^{(\mathrm{e})}(L_{\mathbf{y}}v) = (\mathrm{D}v)(\mathbf{y})\) .
\end{enumerate}

Let \(D_{i}\) be the derivation \(\partial/\partial X_{i}\) of \(A_{x}\) ( \(1 \leqslant i \leqslant n\) ); there exists one and only one left invariant derivation \(T_{i}\) such that \(T_{i}^{(e)} = D_{i}^{(e)}\) . Show that the \(T_{i}\) are linearly independent over K and that every left invariant derivation is a linear combination of the \(T_{i}\) with coefficients in K. Deduce that, for the bracket {[}D, D'{]} and the p-mapping \(D \mapsto D^{p}\) (when p \textgreater{} 0), the set g of left invariant derivations is a Lie algebra (a Lie p-algebra if p \textgreater{} 0) called the Lie algebra of the formal Lie group G.

\begin{enumerate}
\def\labelenumi{(\alph{enumi})}
\setcounter{enumi}{2}
\tightlist
\item
  Show that for every formal power series \(u \in \mathbf{A}_{\mathbf{x}}\) :
\end{enumerate}

\[
u (\mathbf {f} (\mathbf {x}, \mathbf {y})) = u (0) + \sum_ {i = 1} ^ {n} Y _ {i} T _ {i} u + o _ {2} (\mathbf {x}, \mathbf {y})\tag{1}
\]

where in the series \(o_{2}\) all the monomials are of degree \(\geqslant2\) in the \(Y_{j}\) . Deduce that

\[
\begin{array}{r l} {u (\mathbf {f} (\mathbf {x}, \mathbf {f} (\mathbf {y}, \mathbf {z}))) = u (0) + \sum_ {i = 1} ^ {n} (\mathrm{Y} _ {i} + \mathrm{Z} _ {i}) \mathrm{T} _ {i} u} & \\ {+ \sum_ {i, j} \mathrm{Y} _ {i} \mathrm{Z} _ {j} ((\mathrm{T} _ {i} \mathrm{T} _ {j}) (u)) + o _ {3} (\mathbf {x}, \mathbf {y}, \mathbf {z})} & \end{array}\tag{2}
\]

where in the series \(o_{3}\) all the monomials are of degree \(\geqslant3\) in the set of \(Y_{i}\) and \(Z_{i}\) . Deduce that:

\[
u (\mathbf {f} (\mathbf {x}, \mathbf {y})) - u (\mathbf {f} (\mathbf {y}, \mathbf {x})) = \sum_ {i <   j} (X _ {i} Y _ {j} - X _ {j} Y _ {i}) ([ T _ {i}, T _ {j} ] ^ {(e)} (u)) + o _ {3} ^ {\prime} (\mathbf {x}, \mathbf {y})\tag{3}
\]

where all the terms of \(o_3'\) are of total degree \(\geqslant 3\) . Show that, if \(G\) is commutative, \(g\) is commutative.

\begin{enumerate}
\def\labelenumi{(\alph{enumi})}
\setcounter{enumi}{3}
\tightlist
\item
  Let G' be another formal group of dimension m, whose group law is f'. A formal homomorphism of G into G' is a system \(\mathbf{F} = (\mathrm{F}_{j}(\mathbf{x}))_{1 \leqslant j \leqslant m}\) of elements of \(\mathbf{A}_{\mathbf{x}}\) such that \(\mathbf{f}'(\mathbf{F}(\mathbf{x}), \mathbf{F}(\mathbf{y})) = \mathbf{F}(\mathbf{f}(\mathbf{x}, \mathbf{y}))\) . For every left invariant derivation \(D \in g\) , show that \(v \mapsto D(v \circ F)\) is a left invariant derivation belonging to the Lie algebra \(g'\) of G'. If it is denoted by \(\mathbf{F}^*(D)\) , \(\mathbf{F}^*\) is a homomorphism (a p-homomorphism for p \textgreater{} 0) of g into \(g'\) . If \((\mathrm{T}_{1})_{1 \leqslant i \leqslant n}, (\mathrm{T}'_{j})_{1 \leqslant j \leqslant m}\) are the bases of g and \(g'\) such that \(T_{i}^{(e)} = D_{i}^{(e)}\) and \(T'_{j}{}^{(e)} = D'_{j}{}^{(e)}\) , show that the matrix of \(F^*\) relative to these bases is the matrix \((\partial F_{j}/\partial X_{i})_{0}\) (constant term of \(\partial F_{j}/\partial X_{i}\) , where i is the index of the rows and j that of the columns).
\end{enumerate}

\begin{enumerate}
\def\labelenumi{\arabic{enumi}.}
\setcounter{enumi}{24}
\tightlist
\item
  The formal linear group in \(n\) variables over the field \(\mathbf{K}\) , denoted by \(\mathbf{GL}(n, \mathbf{K})\) (when no confusion can arise), is the formal group of dimension \(n^2\) whose group law is defined by \(f_{ij}(\mathbf{u}, \mathbf{v}) = -\delta_{ij} + \sum_{k=1}^{n} (\delta_{ik} + \mathrm{U}_{ik})(\delta_{kj} + \mathrm{V}_{kj})\) ( \(\delta_{ij}\) the Kronecker index; the \(3n^2\) indeterminates of the general definition of Exercise 24 are here denoted by \(\mathrm{U}_{ij}, \mathrm{V}_{ij}, \mathrm{W}_{ij}\) with 2 indices varying from 1 to \(n\) ). Show that, if \(\mathrm{D}_{ij} = \partial/\partial \mathrm{U}_{ij}\) , the left invariant derivations \(\mathrm{X}_{ij}\) such that \(\mathrm{X}_{ij}^{(e)} = \mathrm{D}_{ij}^{(e)}\) are given by
\end{enumerate}

\[
\mathrm{X} _ {i j} = (1 + \mathrm{U} _ {i i}) \mathrm{D} _ {i j} + \sum_ {k \neq i} \mathrm{U} _ {k i} \mathrm{D} _ {k j}.
\]

The Lie algebra of \(\mathbf{GL}(n,\mathbf{K})\) is identified with \(\mathfrak{gl}(n,\mathbf{K})\) by identifying \(X_{ij}\) with the element \(E_{ij}\) of the canonical basis (use formula (3) of Exercise 24). If K is of characteristic p \textgreater{} 0, \(X_{ii}^{p} = X_{ii}\) and \(X_{ij}^{p} = 0\) for \(i \neq j\) .

¶ 26. (a) Let K be a field of characteristic \(\neq 2\) , E an \(n\) -dimensional vector space over K and \(\Phi\) a non-degenerate symmetric bilinear form on E. Suppose that E has an orthogonal basis with respect to \(\Phi\) and let \(R\) denote the (diagonal) matrix of \(\Phi\) with respect to this basis. We know (Algebra, Chapter IX, § 4, Exercise 11) that every matrix (relative to the basis in question ) \(U\) of the orthogonal group \(\mathbf{G}(\Phi)\) such that \(\det(I + U) \neq 0\) can be written

\[
U = (I - R ^ {- 1} S) ^ {- 1} (I + R ^ {- 1} S),
\]

  where \(S = R(U - I)(U + I)^{-1}\) is an alternating matrix such that \(\det(I - R^{-1}S) \neq 0\) ; conversely, for every matrix \(S\) satisfying these conditions, \(U = (I - R^{-1}S)^{-1}(I + R^{-1}S)\) is a matrix of \(\mathbf{G}(\Phi)\) such that \(\det(I + U) \neq 0\) . Let \(S_{ij}, S_{ij}', S_{ij}''\) be \(3n(n - 1)/2\) indeterminates ( \(1 \leqslant i < j \leqslant n\) )

and let L be a field containing the ring of formal power series \(K[[S_{ij}, S'_{ij}, S''_{ij}]]\) . If we denote by \(S, S', S''\) the matrices

\[
\sum_ {i <   j} \mathrm{S} _ {i j} (E _ {i j} - E _ {j i}), \qquad \sum_ {i <   j} \mathrm{S} _ {i j} ^ {\prime} (E _ {i j} - E _ {j i}), \qquad \sum_ {i <   j} \mathrm{S} _ {i j} ^ {\prime \prime} (E _ {i j} - E _ {j i}),
\]

then \(\operatorname{det}(I - R^{-1}S) \neq 0\) and the analogues for \(S'\) and \(S''\) hold; if we write

\[
U = (I - R ^ {- 1} S) ^ {- 1} (I + R ^ {- 1} S), \quad U ^ {\prime} = (I - R ^ {- 1} S ^ {\prime}) ^ {- 1} (I + R ^ {- 1} S ^ {\prime}),
\]

then also \(\operatorname{det}(I + UU') \neq 0\) . We write

\[
F (S, S ^ {\prime}) = R \left(U U ^ {\prime} - I\right) \left(U U ^ {\prime} + I\right) ^ {- 1} = \left(f _ {i j} \left(S, S ^ {\prime}\right)\right);
\]

the \(f_{ij}(S, S')\) belong to the ring of formal power series \(K[[S_{ij}, S_{ij}']]\) and, by considering \(F(S, S')\) as a formal power series with coefficients in the algebra of matrices of order n over K, we can write:

\[
F (S, S ^ {\prime}) = S + S ^ {\prime} + o _ {2} (S, S ^ {\prime})
\]

where in the elements of the matrix \(o_2(S, S')\) the terms are all of degree \(\geqslant 1\) in the \(S_{ij}\) and of degree \(\geqslant 1\) in the \(S_{ij}'\) ; then similarly:

\[
U = I + 2 R ^ {- 1} S + o _ {2} ^ {\prime} (S)
\]

the elements of \(o_2'(S)\) being formal power series of order \(\geqslant 2\) . Show that \(F(S, S')\) is a formal group law over \(K\) of dimension \(n(n - 1)/2\) ; this formal group is called the formal orthogonal group (corresponding to \(\Phi\) ) and is denoted by \(G(\Phi)\) by an abuse of notation.

\begin{enumerate}
\def\labelenumi{(\alph{enumi})}
\setcounter{enumi}{1}
\tightlist
\item
  If we write \(M(S) = (m_{ij}(S)) = (I - R^{-1}S)^{-1}(I + R^{-1}S)\) , \(M\) is a formal homomorphism of \(\mathbf{G}(\Phi)\) into \(\mathbf{GL}(n, \mathbf{K})\) . Let \(\mathfrak{o}(\Phi)\) be the Lie algebra of \(\mathbf{G}(\Phi)\) and let \((\mathrm{T}_{ij})_{i < j}\) denote the basis of this algebra such that \(\mathrm{T}_{ij}^{(e)} = \mathrm{D}_{ij}^{(e)}\) (cf.~Exercise 24 (d)); if an element \(D = \sum_{i < j} c_{ij} T_{ij}\) of \(\mathfrak{o}(\Phi)\) is identified with the matrix \(D = \sum_{i < j} c_{ij}(E_{ij} - E_{ji})\) , show that (in the notation of Exercise 24 (d) with the identification made in Exercise 25 of the Lie algebra of \(\mathbf{GL}(n, \mathbf{K})\) with \(\mathfrak{gl}(n, \mathbf{K})\) ):
\end{enumerate}

\[
M ^ {*} (D) = 2 R ^ {- 1} D.
\]

  Deduce that \(M^*\) is an isomorphism of \(\mathfrak{o}(\Phi)\) onto the subalgebra of \(\mathfrak{gl}(n, \mathbf{K})\) consisting of the matrices \(X\) such that \({}^{t}XR + RX = 0\) (for every ordered pair \((a, b)\) of elements of \(E\) , identified with matrices with one column, consider \(\Phi(a + M(S).a, b + M(S).b) = {}^{t}a.(I + {}^{t}M(S))R(I + M(S)).b\) as a formal power series reduced to its constant term and use the fact that the subalgebra of \(X \in \mathfrak{gl}(n, \mathbf{K})\) such that \({}^{t}XR + RX = 0\) is of dimension \(n(n - 1)/2\) ).

\begin{enumerate}
\def\labelenumi{(\alph{enumi})}
\setcounter{enumi}{2}
\tightlist
\item
  Define similarly the formal symplectic group in 2n variables over an arbitrary
\end{enumerate}

 (commutative) field and show that its Lie algebra is identified with the subalgebra of \(\mathfrak{gl}(2n,\mathbf{K})\) consisting of the matrices X such that

\[
{ } ^ { t } X A + A X = 0 , \quad \text {   where   } \quad A = \left( \begin{array} { c c } 0 & I _ { n } \\ - I _ { n } & 0 \end{array} \right) .
\]

\begin{enumerate}
\def\labelenumi{\arabic{enumi}.}
\setcounter{enumi}{26}
\tightlist
\item
  Let K be a field of characteristic p \textgreater{} 0 and G the 2-dimensional formal group defined by \(f_{1}(\mathbf{x}, \mathbf{y}) = \mathrm{X}_{1} + \mathrm{Y}_{1} + \mathrm{X}_{1}\mathrm{Y}_{1}, f_{2}(\mathbf{x}, \mathbf{y}) = \mathrm{X}_{2} + \mathrm{Y}_{2}(1 + \mathrm{X}_{1}^{p})\) . Show that G is not commutative but that its Lie algebra is commutative.
\end{enumerate}

